\chapter{Đào sâu: Bảo đảm giao nhận Kafka --- ít nhất một lần, consumer lũy đẳng và dead letter}
\label{ch:dd-kafka}

Chương~10 đưa ra quy tắc ngón tay cái --- Kafka là ít-nhất-một-lần, nên
hãy làm consumer lũy đẳng --- và Chương~11 cho thấy phòng thủ poison
pill. Chương đào sâu này đọc các thiết lập producer và consumer mà dự
án thực sự chạy, phần lớn nó chưa từng viết ra, và đi theo một thông
điệp qua mọi nơi nó có thể bị nhân đôi hoặc mất. Rồi bổ sung hai thứ mà
một consumer production cần và dự án còn thiếu: một bản ghi khử trùng
và một đường thử lại kết thúc ở dead-letter topic thay vì một dòng log.

Ba câu hỏi sắp xếp chương này:

\begin{enumerate}
  \item \textbf{Dự án đứng ở đâu} trên thang nhiều-nhất-một-lần /
  ít-nhất-một-lần / đúng-một-lần --- không phải theo ý định, mà theo
  các mặc định nó thừa hưởng?
  \item \textbf{Chính xác khi nào một thông điệp được giao hai lần}, và
  consumer phải làm gì để điều đó vô hại?
  \item \textbf{Chuyện gì xảy ra với một thông điệp không xử lý được},
  và nó nên đi đâu thay vì đi vào hư không?
\end{enumerate}

\section{Vì sao chọn nó, và nó giải quyết gì}

Dự án dùng Kafka cho một việc: tách một lần ghi nghiệp vụ khỏi các hiệu
ứng phụ của nó. Một user đăng ký; một mail phải đi; một dòng học viên
phải xuất hiện trong cơ sở dữ liệu khác. Làm những việc đó inline sẽ
khiến đăng ký chậm và mong manh như SMTP và một Postgres thứ hai. Làm
qua một log khiến chúng \emph{cuối cùng} xảy ra, sống sót qua khởi động
lại consumer, và phát lại được. Chương~6 đã nói về tách đồng bộ/bất
đồng bộ; điều nó chưa nói là ``cuối cùng'' có nghĩa chính xác trong
Kafka, gồm năm hay sáu thiết lập, và mặc định của các thiết lập đó là
thứ dự án đang chạy hôm nay.

\begin{decision}[sự kiện, không phải lệnh --- và khi nào một lời gọi mới đúng]
\code{ts.user.registered} nói ``việc này đã xảy ra''; nó không nói
``hãy gửi e-mail''. Notification-service tự quyết đăng ký nghĩa là gì
với nó, và course-service bỏ qua topic này hoàn toàn. Đó là phong cách
sự kiện, và là lý do một consumer thứ tư có thể thêm vào mà không đụng
auth-service. Phong cách lệnh --- một thông điệp \code{SendEmail} gửi
đích danh cho một handler --- ghép bên gửi vào sự tồn tại của bên nhận
và thường là dấu hiệu rằng người ta thực ra muốn một lời gọi đồng bộ.
Và đôi khi đúng là thế: khi bên gọi cần \emph{kết quả} ngay (user này có
tồn tại không? file này đã quét chưa?), một lời gọi Feign với circuit
breaker là thiết kế trung thực, và khoác cho nó cái áo request-reply
qua Kafka chỉ thêm hai chặng và một correlation id. Quy tắc của dự án:
sự kiện cho các sự thật mà service khác phản ứng; lời gọi cho các câu
hỏi bên gọi phải có đáp án trước khi tiếp tục.
\end{decision}

\section{Những thiết lập dự án chưa từng viết}

\subsection{Producer}

Course-service dựng producer bằng tay trong \code{KafkaProducerConfig}:

\begin{lstlisting}[language=Java]
Map<String, Object> configProps = new HashMap<>();
configProps.put(ProducerConfig.BOOTSTRAP_SERVERS_CONFIG, bootstrapServers);
configProps.put(ProducerConfig.KEY_SERIALIZER_CLASS_CONFIG,
        StringSerializer.class);                              (*@\co{1}@*)
configProps.put(ProducerConfig.VALUE_SERIALIZER_CLASS_CONFIG,
        JsonSerializer.class);
return new DefaultKafkaProducerFactory<>(configProps);
\end{lstlisting}

Ba dòng --- và mọi thứ khác là mặc định của client Kafka~3.9. Từ
Kafka~3.0 các mặc định đó là những giá trị an toàn:

\begin{lstlisting}[language=properties]
acks=all
enable.idempotence=true
retries=2147483647
max.in.flight.requests.per.connection=5
delivery.timeout.ms=120000
max.block.ms=60000
linger.ms=0
\end{lstlisting}

Auth-service và notification-service dùng factory tự cấu hình của
Spring Boot từ YAML, thứ đặt cùng hai serializer và không gì khác, nên
cả ba producer hành xử giống nhau.

\begin{description}
  \item[\co{1}] Key serializer quan trọng hơn vẻ ngoài. Mọi publisher
  trong dự án đều truyền key: \code{event.userId()} trong
  \code{UserRegisteredPublisher} và \code{UserActivatedPublisher},
  \code{assignmentId} trong \code{AssignmentCreatedPublisher}. Kafka
  băm key để chọn partition; key bằng nhau, partition bằng nhau, và
  broker giữ thứ tự trong một partition. Đó là toàn bộ bảo đảm thứ tự
  theo user, và nó là một chuỗi ký tự.
\end{description}

Các mặc định nghĩa là gì, từng cái một:

\begin{itemize}
  \item \textbf{\code{acks=all}}: leader chờ mọi replica đồng bộ trước
  khi xác nhận. Với \code{replicas: 1} trong \code{kafka-topics.yaml} và
  một broker duy nhất trong Compose, ``all'' là một --- độ bền bằng một
  cái đĩa. Trong production nó đi cùng \code{replication.factor=3} và
  \code{min.insync.replicas=2}: một lần ghi sống sót khi một broker
  chết, và nếu hai broker chết producer nhận \code{NotEnoughReplicas}
  thay vì một lần ghi bản-sao-đơn âm thầm.
  \item \textbf{Producer lũy đẳng}: mỗi producer nhận một id và đánh số
  thông điệp theo partition; một lần thử lại sau khi mất ack được broker
  nhận ra và không nối thêm lần hai. Điều này khép lại bản sao \emph{phía
  producer} của Chương~10. Nó không làm gì cho phía consumer.
  \item \textbf{\code{retries} không giới hạn, bị chặn bởi
  \code{delivery.timeout.ms}}: producer thử lại lỗi tạm thời tới hai
  phút, rồi làm future thất bại. \code{send()} đã trả về từ lâu; không
  ai trong dự án đọc future đó.
\end{itemize}

\begin{pitfall}[YAML không làm gì cả]
Triệu chứng: bạn thêm \code{spring.kafka.producer.acks: 1} vào
\code{course-service.yml} để tăng tốc một bài load test, và không gì
thay đổi. Nguyên nhân: \code{KafkaProducerConfig} khai báo bean
\code{ProducerFactory} riêng, việc đó tắt auto-configuration của Boot
cho factory ấy --- khối producer trong YAML không được ai đọc trong
course-service. Auth-service, không có bean như vậy, \emph{sẽ} tuân
theo nó. Hai service, một quy ước YAML, hai hành vi. Cách sửa: xóa bean
viết tay (factory của Boot làm y hệt) hoặc chuyển mọi thuộc tính vào
đó; không bao giờ cả hai.
\end{pitfall}

\subsection{Consumer}

YAML consumer trong \code{notification-service.yml} và
\code{course-service.yml} đặt tên group, \code{auto-offset-reset:
earliest}, và các deserializer xử lý lỗi. Mọi thứ về \emph{commit} lại
là mặc định, lần này của Spring Kafka:

\begin{itemize}
  \item \code{enable.auto.commit=false}. Spring tắt commit theo hẹn giờ
  của client và tự commit.
  \item \code{AckMode.BATCH}. Sau khi listener trả về cho \emph{mọi}
  record của một \code{poll()}, Spring commit offset cao nhất của lô.
  Xử lý và commit cách nhau bằng thời gian xử lý của lô.
  \item \code{DefaultErrorHandler} với \code{FixedBackOff(0, 9)}: một
  listener ném ngoại lệ được thử lại thêm chín lần không chờ, rồi
  record được \emph{log và bỏ qua}. Không gửi đi đâu. Không đếm ở đâu
  mà dự án đọc.
  \item \code{max.poll.records=500}, \code{max.poll.interval.ms=300000}:
  một lô tới 500 record phải xử lý xong trong năm phút, nếu không
  consumer bị coi là đã chết.
\end{itemize}

Ghép bốn thứ đó lại và vị trí của dự án là: ít-nhất-một-lần, với bản
sao khả dĩ ở mọi lần crash hay rebalance, và với mất mát âm thầm sau
mười lần thử thất bại. Cả thứ tự lẫn mất mát đều không được viết trong
file nào bạn có thể \code{grep}.

\section{Chính xác khi nào xảy ra bản sao}

Dòng thời gian tạo ra một bản sao, với \code{notification-group} chạy
hai instance A và B trên một topic có hai partition:

\begin{lstlisting}
t0   A polls partition 0, gets offsets 40..49   (committed: 40)
t1   A processes 40..47: eight e-mails sent
t2   A is killed (rolling update, OOMKilled, SIGKILL)
     -- offset 40 is still the committed position --
t3   group rebalances; partition 0 is assigned to B
t4   B polls from the committed offset: 40
t5   B processes 40..49: ten e-mails sent
     -- users 40..47 receive their activation mail twice --
\end{lstlisting}

Không gì ở đây là lỗi. Nó là hợp đồng: Kafka hứa offset sẽ không tiến
qua một record bạn chưa xử lý xong; cái giá là một record bạn \emph{đã}
xử lý xong nhưng chưa commit sẽ quay lại. Mọi consumer làm bất cứ điều
gì không đảo ngược được --- gửi mail, chèn dòng, trừ tiền thẻ --- phải
được viết cho \code{t5}.

Hai biến thể quan trọng với dự án này. \code{AckMode.RECORD} commit sau
mỗi record thay vì mỗi lô: ít bản sao hơn khi crash (một thay vì tới
500), nhiều lượt commit hơn. Và biến thể \emph{consumer chậm} không cần
crash nào: nếu xử lý một lô lâu hơn \code{max.poll.interval.ms}, broker
loại consumer khỏi group, instance khác nhận partition từ commit cuối,
và \emph{cả hai} cuối cùng xử lý cùng những record --- cái đầu chưa
nhận ra mình đã bị đuổi.

\begin{pitfall}[một bài tập, một cơn bão e-mail trùng]
Triệu chứng: giáo viên tạo bài tập trong khóa có 300 học viên; một số
học viên nhận thông báo ba lần; log consumer hiện \code{Member ... has
left the group} và một rebalance. Nguyên nhân:
\code{AssignmentCreatedListener} lặp qua \code{enrolledStudentEmails} và
gọi \code{javaMailSender.send()} cho từng địa chỉ, đồng bộ, trên luồng
poll. Một giây mỗi vòng SMTP, 300 e-mail là năm phút --- đúng bằng
\code{max.poll.interval.ms}. Consumer bị đuổi giữa vòng lặp, partition
được gán lại, và chủ mới bắt đầu vòng lặp lại từ học viên đầu tiên.
Cách sửa, theo chi phí: tăng \code{max.poll.interval.ms} (vá tạm), phát
một record \code{ts.notification.email} cho mỗi người nhận để mỗi mail
là một đơn vị công việc riêng (hình dạng đúng), và làm việc gửi lũy
đẳng theo người nhận (Mục~\ref{sec:enhancing-kafka}).
\end{pitfall}

\section{Ưu và nhược, thẳng thắn}

\begin{center}
\begin{tabular}{@{}p{0.28\linewidth}p{0.32\linewidth}p{0.32\linewidth}@{}}
\toprule
 & Dự án hôm nay & Nền production \\
\midrule
Độ bền producer & \code{acks=all}, 1 replica & 3 replica, \code{min.insync}=2 \\
Bản sao producer & không (lũy đẳng) & không (lũy đẳng) \\
Bản sao consumer & có thể, chưa xử lý & có thể, đã khử trùng \\
Record thất bại & log và bỏ sau 10 lần & topic thử lại, rồi DLT \\
Thứ tự & theo key, 1 partition & theo key, N partition \\
Việc chậm & trên luồng poll & tách ra, hoặc bất đồng bộ \\
Quan sát & không có lag & cảnh báo lag, cảnh báo DLT \\
\bottomrule
\end{tabular}
\end{center}

Điểm yếu trung thực của dự án: mất mát âm thầm sau mười lần thử; fan-out
e-mail trên luồng poll; không khử trùng nơi hiệu ứng phụ không tự nhiên
lũy đẳng (mail); không có metric lag nào ai nhìn; và topic một partition,
khiến thứ tự tầm thường và song song bất khả.

\section{Nâng cấp giải pháp}
\label{sec:enhancing-kafka}

\subsection{Một consumer lũy đẳng có sổ ghi những gì đã thấy}

\code{UserActivatedListener} đã lũy đẳng theo khóa tự nhiên --- giao
lại lần hai tìm thấy học viên và trả về. Nơi không có khóa tự nhiên, ghi
lại id sự kiện. Outbox của Chương~33 đặt một header \code{eventId} duy
nhất lên mọi thông điệp; consumer giữ một bảng các id nó đã xử lý xong,
và ghi dòng đó \emph{trong cùng transaction} với thay đổi nghiệp vụ.
Flyway \code{V3\_\_create\_processed\_events.sql} trong course-service:

\begin{lstlisting}[language=SQL]
CREATE TABLE processed_events (
    event_id     UUID PRIMARY KEY,
    consumer     VARCHAR(64) NOT NULL,     -- listener name
    processed_at TIMESTAMPTZ NOT NULL DEFAULT now()
);
-- retention: rows older than the topic's retention are useless
CREATE INDEX idx_processed_events_at ON processed_events (processed_at);
\end{lstlisting}

Listener, viết lại để đọc header và giành id trước:

\begin{lstlisting}[language=Java]
@KafkaListener(topics = "ts.user.activated", groupId = "course-group")
@Transactional                                              (*@\co{1}@*)
public void onUserActivated(
        UserActivatedEvent event,
        @Header("eventId") String eventId) {

    int claimed = jdbc.update("""
        INSERT INTO processed_events (event_id, consumer)
        VALUES (?::uuid, 'UserActivatedListener')
        ON CONFLICT (event_id) DO NOTHING""", eventId);      (*@\co{2}@*)
    if (claimed == 0) {
        log.info("duplicate event {} ignored", eventId);
        return;                                             (*@\co{3}@*)
    }
    if (!ROLE_STUDENT.equalsIgnoreCase(event.role())) {
        return;
    }
    Long userId = Long.parseLong(event.userId());
    Student student = studentRepository.findByUserId(userId)
            .orElseGet(() -> Student.builder().userId(userId).build());
    student.setFirstName(event.firstName());               (*@\co{4}@*)
    student.setLastName(event.lastName());
    student.setEmail(event.email());
    studentRepository.save(student);
}
\end{lstlisting}

\begin{description}
  \item[\co{1}] Một transaction cơ sở dữ liệu bao quanh việc giành và
  việc thay đổi. Nếu \code{save()} thất bại, lượt giành cũng rollback,
  và lần thử lại có thể giành lại. Nếu JVM chết sau commit nhưng trước
  khi commit offset Kafka, lần giao lại rơi vào \co{3}.
  \item[\co{2}] \code{ON CONFLICT DO NOTHING} là toàn bộ thuật toán khử
  trùng: một insert-nếu-chưa-có nguyên tử. Không có race đọc-rồi-ghi
  giữa hai instance consumer cùng nhận một record sau rebalance.
  \item[\co{3}] Trả về bình thường, không ném ngoại lệ: offset phải tiến
  qua một bản sao.
  \item[\co{4}] Nhân tiện, upsert thay vì insert-nếu-thiếu, để cùng
  listener này phục vụ luôn sự kiện \code{ts.user.updated} mà Chương~33
  thêm vào.
\end{description}

Notification-service không có cơ sở dữ liệu, và mail là hiệu ứng phụ
duy nhất bạn không thể thu hồi. Các lựa chọn, thẳng thắn về chi phí:
chấp nhận thỉnh thoảng một mail trùng (điều dự án ngầm làm); hoặc cho
service một kho nhỏ --- Redis \code{SET eventId 1 NX EX 604800} trước
khi gửi, cùng lượt giành đó trong một lệnh; hoặc làm chính mail lũy đẳng
bằng cách đặt id sự kiện vào header \code{Message-ID}, để hộp thư khử
trùng thứ consumer không làm được.

\subsection{Topic thử lại và dead-letter topic}

Thay ``thử lại mười lần trong một mili giây, rồi quên'' bằng thử lại có
chờ, và một đích đến cho thứ vẫn thất bại. \code{@RetryableTopic} của
Spring Kafka làm điều này bằng topic thay vì thread, nên một record kẹt
không chặn partition phía sau nó:

\begin{lstlisting}[language=Java]
@RetryableTopic(
        attempts = "4",                                     (*@\co{1}@*)
        backoff = @Backoff(delay = 2_000, multiplier = 3),  (*@\co{2}@*)
        include = { MailSendException.class,
                    MailAuthenticationException.class },    (*@\co{3}@*)
        dltStrategy = DltStrategy.FAIL_ON_ERROR)
@KafkaListener(topics = "ts.user.registered",
               groupId = "notification-group")
public void onUserRegistered(UserRegisteredEvent event) {
    emailService.sendActivationEmail(
            event.email(), event.firstName(), event.activationCode());
}

@DltHandler                                                 (*@\co{4}@*)
public void onDeadLetter(UserRegisteredEvent event,
        @Header(KafkaHeaders.EXCEPTION_MESSAGE) String reason,
        @Header(KafkaHeaders.ORIGINAL_OFFSET) long offset) {
    log.error("DLT ts.user.registered offset={} user={} reason={}",
            offset, event.userId(), reason);
    meter.counter("kafka.dlt", "topic", "ts.user.registered")
         .increment();                                      (*@\co{5}@*)
}
\end{lstlisting}

\begin{description}
  \item[\co{1}] Một lần thử trên topic chính, ba lần trên các topic thử
  lại.
  \item[\co{2}] Độ trễ 2, 6 và 18 giây --- khoảng nửa phút để MailDev
  hay một SMTP relay quay lại. Spring tạo các topic
  \code{ts.user.registered-retry-2000}, \code{-retry-6000},
  \code{-retry-18000} và \code{ts.user.registered-dlt}; dưới Strimzi
  (Chương~25) hãy khai báo chúng là object \code{KafkaTopic} hoặc để
  auto-creation bật.
  \item[\co{3}] Chỉ thử lại thứ mà chờ đợi sửa được. Lỗi giải tuần tự
  hay \code{NullPointerException} đi thẳng tới DLT; thử lại nó ba lần là
  diễn kịch.
  \item[\co{4}] Handler dead-letter là một listener bình thường trên
  topic \code{-dlt}. Nó \emph{không được} ném ngoại lệ.
  \item[\co{5}] Một counter, để ``DLT không rỗng'' là một cảnh báo (khả
  năng quan sát của Chương~27) thay vì một phát hiện tình cờ.
\end{description}

Làm gì với một dead letter là quyết định vận hành, không phải code: một
cảnh báo bắn; ai đó đọc header \code{EXCEPTION\_MESSAGE}; bản sửa được
đưa lên; và các record được \emph{phát lại} --- chép từ DLT về topic
chính bằng một công cụ admin nhỏ hoặc Kafka UI --- nơi consumer lũy đẳng
bảo đảm những cái đã thành công bị bỏ qua.

\begin{handson}[xem bậc thang thử lại và dead letter]
Làm hỏng SMTP, đăng ký một user, và đọc các topic:
\begin{lstlisting}[language=bash]
$ docker compose stop maildev
$ curl -s -X POST localhost:8080/api/auth/register -H \
  'Content-Type: application/json' -d '{"email":"dlt@test.io", \
  "password":"Passw0rd!","firstName":"D","lastName":"L", \
  "activationMethod":"EMAIL"}'
$ docker compose exec kafka /opt/kafka/bin/kafka-topics.sh \
  --bootstrap-server localhost:19092 --list | grep registered
ts.user.registered
ts.user.registered-dlt
ts.user.registered-retry-18000
ts.user.registered-retry-2000
ts.user.registered-retry-6000
\end{lstlisting}
Khoảng 26 giây sau request, notification-service log
\code{DLT ts.user.registered offset=0 user=... reason=Mail server
connection failed}. Khởi động lại MailDev, chép record từ DLT sang
\code{ts.user.registered} trong Kafka UI (\code{localhost:9090}), và
mail kích hoạt tới nơi.
\end{handson}

\subsection{Đúng một lần, và những gì nó không bao phủ}

Transaction của Kafka cho đúng-một-lần với một hình dạng cụ thể: tiêu
thụ từ topic A, sản xuất sang topic B, commit cả offset lẫn đầu ra một
cách nguyên tử (\code{isolation.level=read\_committed} ở bên đọc). Các
bộ xử lý luồng cần nó; dự án không có giai đoạn nào như vậy. Transaction
\emph{không} trải qua Kafka và Postgres, hay Kafka và SMTP --- khoảnh
khắc hiệu ứng phụ rời khỏi broker, bạn quay về ít-nhất-một-lần cộng lũy
đẳng. Người phỏng vấn hỏi ``có thể đạt đúng-một-lần không?'' muốn đúng
sự phân biệt đó: có bên trong Kafka, không xuyên hệ thống, và outbox
cộng một bảng khử trùng là cách bạn đạt hiệu quả đó xuyên hệ thống dù
sao đi nữa.

\section{Chuyện gì gãy lúc 3 giờ sáng}

\begin{itemize}
  \item \textbf{notification-service chết bốn giờ.} Không gì mất: record
  tích lại trên topic, retention là 168 giờ trong Compose và mặc định
  bảy ngày của Strimzi. Consumer lag --- khoảng cách giữa offset cuối
  của partition và offset đã commit của group --- lớn lên vài nghìn.
  Khi khởi động lại service tiêu thụ chúng ở tốc độ SMTP; mọi mail kích
  hoạt đến trễ, không cái nào thiếu. Điều \emph{sẽ} làm mất chúng: một
  sự cố dài hơn retention, hoặc một lần triển khai lại dưới group id
  \emph{mới}, bắt đầu từ \code{earliest} và gửi lại mọi mail trong bảy
  ngày lịch sử.
  \item \textbf{Broker chết.} Producer chặn trong \code{max.block.ms}
  ở lần lấy metadata đầu, rồi ném ngoại lệ --- sự cố đăng ký sáu mươi
  giây của Chương~33. Producer đã kết nối sẵn đệm tới
  \code{buffer.memory} (32\,MB) và thử lại trong
  \code{delivery.timeout.ms}; sau hai phút future thất bại và, vì không
  ai đọc nó, sự kiện biến mất. Consumer đơn giản là poll không được gì.
  Outbox biến tất cả chuyện này thành một cột \code{attempts} tăng dần.
  \item \textbf{Phát lại từ hôm qua.} Vận hành reset \code{course-group}
  về một mốc thời gian bằng
  \code{kafka-consumer-groups.sh --reset-offsets --to-datetime}. Mọi sự
  kiện từ đó được giao lại; với lượt giành \code{processed\_events} mỗi
  cái tốn một insert thất bại. Không có nó, mail đi hai lần và dòng học
  viên được thử tạo hai lần (chỉ được cứu bởi ràng buộc
  \code{UNIQUE(user\_id)}).
  \item \textbf{Một key nóng.} Mọi sự kiện cho một khóa học rất lớn rơi
  vào một partition; thêm partition không giúp gì vì key chọn partition.
  Đổi key --- theo người nhận thay vì theo bài tập --- hoặc chấp nhận
  thông lượng của một key là thông lượng của một consumer.
\end{itemize}

\section{Ngoài dự án}

Mọi Kafka được quản lý --- Confluent Cloud, AWS MSK, Aiven --- đều giao
\code{acks=all}, ba replica và \code{min.insync.replicas=2} làm điểm
khởi đầu, và bán dashboard lag như thứ đầu tiên bạn nhìn. Topic thử lại
và DLT chuẩn tới mức Parallel Consumer của Confluent, consumer proxy
của Uber và \code{@RetryableTopic} của Spring đều cài cùng bậc thang;
khác biệt ở cách chúng giữ thứ tự theo key qua các lần thử lại (Spring
mặc định không giữ --- một record thử lại có thể được xử lý sau một
record muộn hơn cùng key, ổn cho mail và không ổn cho bút toán sổ cái).
RabbitMQ và SQS cho bạn ack theo từng thông điệp và dead-letter queue
như tính năng của broker thay vì mẫu thư viện; thứ chúng không thể cho
là phát lại, chính là lý do một log được chọn ngay từ đầu.

\section{Câu hỏi từ phỏng vấn thật}

\subsection*{Q: Làm sao tăng số thông điệp Kafka xử lý được trong một đơn vị thời gian?}

Tôi bắt đầu từ consumer, vì đó là nơi giới hạn của dự án nằm. Partition
là đơn vị song song: một group không thể có nhiều consumer hoạt động
hơn số partition của topic, và mọi topic trong \code{kafka-topics.yaml}
đều có một. Nên thay đổi đầu tiên là số partition --- 6 hoặc 12, chọn
dư vì partition thêm được nhưng không bao giờ bớt được, và thêm sau này
làm đổi phép băm key-sang-partition, phá thứ tự theo key với dữ liệu
đang bay. Rồi khớp consumer với partition:
\code{spring.kafka.listener.concurrency=3} mỗi pod nhân hai replica cho
sáu thread trên sáu partition. Kế đến, đưa việc chậm ra khỏi luồng
poll: một \code{AssignmentCreated} gửi 300 e-mail đồng bộ là một lô năm
phút làm sập \code{max.poll.interval.ms}; thay vào đó phát một record
cho mỗi người nhận, để thông lượng tăng theo partition. Phía producer,
gom lô: \code{linger.ms=5--20} và \code{batch.size=64--256\,KB} cho
client gửi hàng trăm record mỗi request, và \code{compression.type=lz4}
hoặc \code{zstd} cắt byte trên đường truyền ba đến năm lần với JSON.
\code{acks=1} tăng thông lượng producer chừng gấp đôi với cái giá mất
một lần ghi khi leader chết trước khi nhân bản --- tôi giữ
\code{acks=all} cho bất cứ thứ gì không tái tạo được. Phía consumer,
\code{max.poll.records} và \code{fetch.min.bytes} đổi độ trễ lấy ít vòng
đi hơn, bị chặn bởi thứ vừa trong \code{max.poll.interval.ms}. Và tôi để
ý key nóng: nếu một partition gánh 80\,\% tải vì chọn key dở, không gì ở
trên giúp được và cách sửa là một key khác.

\subsection*{Q: Consumer lag đang tăng. Bạn chẩn đoán thế nào?}

Trước hết tôi xác định nó là một trong ba thứ: consumer chậm, consumer
chết, hay producer nhanh lên. \code{kafka-consumer-groups.sh --describe
--group notification-group} cho lag theo partition và member nào sở hữu
mỗi cái; lag trên \emph{mọi} partition mà member vẫn có mặt nghĩa là xử
lý chậm; lag trên \emph{một số} với member biến mất nghĩa là một
instance chết hoặc kẹt trong rebalance; lag xuất hiện ở một thời điểm rõ
ràng với tiêu thụ đều nghĩa là một đợt bùng nổ producer. Với xử lý chậm
tôi nhìn thời gian mỗi record của listener --- trong dự án này là độ trễ
SMTP --- và sự khuấy đảo rebalance trong log (\code{Member ... has
left}), nghĩa là \code{max.poll.interval.ms} đang bị chạm và group đang
quay cuồng. Cách sửa đi theo chẩn đoán: thêm partition và concurrency,
đưa lời gọi chậm ra khỏi luồng, hoặc, nếu đợt bùng nổ là hợp lệ, để nó
chảy hết và chỉ cảnh báo khi lag già hơn SLO tính bằng giây, không phải
lớn hơn một số record. Metric để vẽ là lag theo \emph{thời gian}, không
theo số lượng.

\subsection*{Q: Một thông điệp bị xử lý hai lần. Vì sao, và làm sao cho an toàn?}

Vì Kafka commit offset sau khi xử lý, và tiến trình chết, hoặc bị đuổi
vì chậm, giữa hai việc đó. Đó là ít-nhất-một-lần hoạt động đúng thiết
kế. Làm cho an toàn nghĩa là mọi consumer hoặc có khóa lũy đẳng tự
nhiên --- \code{findByUserId} của \code{UserActivatedListener} --- hoặc
giành một id sự kiện duy nhất trong cùng transaction với hiệu ứng phụ:
\code{INSERT ... ON CONFLICT DO NOTHING} vào \code{processed\_events},
trả về nếu không dòng nào. Id phải tới từ producer, đúc một lần và mang
qua các lần thử lại, chính là thứ dòng outbox cho bạn. Với hiệu ứng phụ
không có transaction --- e-mail --- tôi đặt id vào \code{Message-ID} của
thư hoặc giành nó trong Redis bằng \code{SET NX}. Điều tôi không làm là
chuyển sang \code{AckMode.RECORD} và gọi là đã sửa; việc đó thu hẹp cửa
sổ từ một lô xuống một record và để nguyên lỗi tại chỗ.

\subsection*{Q: Làm sao bảo đảm thứ tự cho một user?}

Key mọi sự kiện về user đó theo id user, để phép băm đặt chúng vào một
partition, và giữ một thread consumer cho mỗi partition. Đó là toàn bộ
bảo đảm, và nó có ba cách vỡ: producer lũy đẳng bị tắt với
\code{max.in.flight} trên một, cho phép một lô thử lại rơi sau lô muộn
hơn; đổi số partition, làm băm lại key; và topic thử lại, chèn lại một
record thất bại phía sau những record kế tiếp nó. Với dự án, key
\code{userId} trên \code{ts.user.*} đã làm việc này. Nếu nghiệp vụ cần
phát biểu mạnh hơn --- ``cập nhật được áp dụng sau tạo, luôn luôn'' ---
tôi thêm số phiên bản vào sự kiện và để consumer từ chối bất cứ thứ gì
cũ hơn cái nó có, điều biến thứ tự từ thuộc tính của broker thành bất
biến của consumer.

\subsection*{Q: Một broker chết. Producer làm gì?}

Với ba broker và \code{replication.factor=3}, controller bầu leader mới
từ các replica đồng bộ cho mỗi partition bị ảnh hưởng trong vài giây;
producer nhận \code{NotLeaderOrFollower}, làm mới metadata, và thử lại
--- vô hình với ứng dụng vì \code{retries} không giới hạn và
\code{delivery.timeout.ms} là hai phút. Nếu broker chết kéo ISR xuống
dưới \code{min.insync.replicas}, producer nhận \code{NotEnoughReplicas}
và tiếp tục thử lại; với \code{acks=all} đó là sự từ chối đúng đắn việc
ghi một bản sao duy nhất. Trên stack Compose một broker của dự án không
có failover: các lần gửi được đệm, future thất bại sau hai phút, và dòng
outbox là thứ cứu sự kiện. Điều duy nhất tôi kiểm tra khi review là có
ai đọc future không --- trong dự án này không ai, nên lỗi phía producer
là âm thầm.

\subsection*{Q: Làm sao phát lại một topic từ hôm qua?}

Dừng consumer group, reset offset của nó, khởi động lại:
\code{kafka-consumer-groups.sh --reset-offsets --group course-group
--topic ts.user.activated --to-datetime 2026-09-06T00:00:00.000
--execute}. Group phải không hoạt động, nếu không lệnh từ chối. Trước
khi chạy tôi trả lời hai câu hỏi: consumer có lũy đẳng với mọi thứ kể
từ mốc đó không --- vì mọi record sẽ quay lại --- và dữ liệu còn đó
không, phụ thuộc vào retention (\code{168h} ở đây) và topic có compact
không. Với một record đơn lẻ tôi phát lại từ DLT thay vì reset cả group.
Và tôi không bao giờ phát lại vào một group gửi e-mail mà chưa xác nhận
đường khử trùng, vì ``chúng ta gửi lại 40\,000 mail kích hoạt'' là sự cố
theo sau một lần reset tự tin.

\begin{quiz}
\begin{enumerate}
  \item Nêu tên các thiết lập producer khiến producer của dự án lũy
  đẳng và bền vững, và giải thích vì sao \code{acks=all} trên stack
  Compose yếu hơn nghe có vẻ.
  \item Vẽ (bằng lời) dòng thời gian trong đó \code{notification-group}
  gửi cùng một mail kích hoạt hai lần mà không có crash nào, và nêu tên
  thiết lập mà giá trị của nó định nghĩa điểm kích hoạt.
  \item Vì sao insert vào \code{processed\_events} và lần ghi nghiệp vụ
  phải chung một transaction? Mô tả sự xen kẽ đi sai nếu insert commit
  trước một mình.
  \item Trong cấu hình \code{@RetryableTopic}, giải thích vì sao
  \code{include} chỉ liệt kê ngoại lệ mail và chuyện gì xảy ra với một
  \code{NullPointerException}.
  \item Course-service khai báo bean \code{ProducerFactory} riêng. Nêu
  hệ quả cho khối \code{spring.kafka.producer} trong YAML của nó, và
  cách sửa một dòng.
  \item Notification-service đã chết chín ngày. Nêu chuyện gì xảy ra
  khi khởi động lại, thiết lập nào quyết định, và điều gì khác đi nếu nó
  khởi động lại với \code{group-id: notification-v2}.
\end{enumerate}
\end{quiz}
